% =========================
% PREAMBLE (preamble.tex)
% =========================

% --- Dil ve kodlama: bunlar en başta olsun ---
\usepackage[T1]{fontenc}
\usepackage[utf8]{inputenc}
\usepackage[turkish,shorthands=off]{babel}


% --- Özet (book sınıfı için) ---
\newcommand{\abstractname}{Özet}
\newenvironment{abstract}{
  \cleardoublepage
  \thispagestyle{plain}
  \chapter*{\abstractname}
  \addcontentsline{toc}{chapter}{\abstractname}
}{\cleardoublepage}


% --- Yazı tipi ve mikro tipografi (önerilir) ---
\usepackage{lmodern}
%\let\CheckCommand\providecommand
%\usepackage{microtype}

% --- Sayfa yerleşimi ---
\usepackage{geometry}
\geometry{margin=2.5cm}

% --- Matematik ve semboller ---
\usepackage{amsmath, amssymb, bm}

% --- Grafik, çizim, tablo ---
\usepackage{graphicx}
\usepackage{float}
\usepackage{tabularx}
\usepackage{longtable}
\usepackage{array}
\usepackage{booktabs}
\usepackage{tikz}
\usepackage{xcolor}
\usetikzlibrary{decorations.pathmorphing}
\usetikzlibrary{arrows.meta,positioning,shapes.geometric,fit,calc}
\usepackage{fvextra}
\usepackage{dirtree}


% --- Kutular / renkli çerçeveler ---
% not: 'tcolor' değil 'tcolorbox' olmalı
\usepackage[most]{tcolorbox}

% --- Diğer yardımcılar ---
\usepackage[normalem]{ulem}
\usepackage{comment}

\usepackage[
  backend=biber,
  style=ieee,
  maxbibnames=3,
  maxcitenames=2,
  giveninits=true,
  doi=false,
  url=false,
  isbn=false
]{biblatex}

\addbibresource{references.bib}

\usepackage{caption}
\usepackage{algorithm}
\usepackage{algorithmic}

% --- Hiperbağlantılar (en sonda yüklensin) ---
\usepackage{hyperref}
\hypersetup{
    unicode=true,           % Türkçe karakterler için
    colorlinks=true,
    linkcolor=black,
    filecolor=black,
    urlcolor=black,
    citecolor=black,
    pdftitle={Başlık},
    pdfauthor={Yazar}
}

% preamble.tex içinde
\usepackage{graphicx}
\usepackage{eso-pic}

\newcommand{\CoverBackground}{%
  \AddToShipoutPictureBG*{%
    \AtPageLowerLeft{%
      \includegraphics[width=\paperwidth,height=\paperheight]{Resimler/Yerçekimi Modelleri.png}%  % Kapa Tasarımı.png dosyasını KapaTasarimi.png yap
    }%
  }%
}

\newcommand{\volume}{{\ooalign{\hfil$V$\hfil\cr\kern0.08em--\hfil\cr}}}